\section{Command line reference}
\label{sec:cli}

\begin{lstlisting}[style=shell]
rubiks-cfop [--scramble MOVES | --facelets STATE | --random]
            [--scramble-style {random-state,random-move}] [--seed N]
            [--cross {best,white,red,green,yellow,orange,blue}]
            [--color | --no-color] [--json] [--cross-options]
\end{lstlisting}

\subsection{Where the cube comes from}

Exactly one of these three sources is used.  With none of them the program behaves
as if \texttt{--random} had been given.

\begin{longtable}{@{}p{3.2cm}p{11.2cm}@{}}
\toprule
Option & Meaning \\
\midrule
\endhead
\texttt{--scramble MOVES} &
Apply a scramble in WCA notation to a solved cube.  Wide moves
(\texttt{Rw}, \texttt{2Rw}), rotations (\texttt{x y z}) and slice moves
(\texttt{M E S}) are accepted, as are the usual aliases (\texttt{r} for
\texttt{Rw}, \texttt{m} for \texttt{M}).  Invalid notation is reported with the
offending token and its position, and the program exits with status 2. \\
\texttt{--facelets STATE} &
Start from a 54 character Kociemba facelet string instead of a scramble.  Useful
when another tool produced the state, or when you want to solve a state you built
by hand.  Whitespace and lower case are ignored; impossible states (a twisted
corner, a wrong colour count, moved centres) are rejected. \\
\texttt{--random} &
Generate a scramble and solve it.  This is the default. \\
\bottomrule
\end{longtable}

\subsection{How the scramble is generated}

\begin{longtable}{@{}p{3.2cm}p{11.2cm}@{}}
\toprule
Option & Meaning \\
\midrule
\endhead
\texttt{--scramble-style} &
\texttt{random-state} (default) draws a uniformly random cube state and solves it,
which is what WCA Regulation 4b3 requires; scrambles are 20 to 22 moves.
\texttt{random-move} produces a filtered random-move sequence of 20 face turns
instead --- instant, but \emph{not} WCA-legal for the $3\times3\times3$. \\
\texttt{--seed N} &
Seed for a generated scramble.  The same seed always gives the same scramble, and
the seed is echoed in the JSON payload.  Without a seed one is drawn from the
operating system and reported, so any run can be reproduced. \\
\bottomrule
\end{longtable}

\subsection{How it is solved}

\begin{longtable}{@{}p{3.2cm}p{11.2cm}@{}}
\toprule
Option & Meaning \\
\midrule
\endhead
\texttt{--cross} &
The colour of the cross: \texttt{white}, \texttt{yellow}, \texttt{green},
\texttt{blue}, \texttt{red} or \texttt{orange}.  \texttt{best} (the default) solves
all six crosses and keeps the shortest, breaking ties in the order white, yellow,
green, blue, red, orange.  The cube is always reoriented so that the chosen colour
ends up in the \emph{bottom} layer. \\
\texttt{--cross-options} &
Before solving, print the six crosses with their lengths, shortest first. \\
\bottomrule
\end{longtable}

\subsection{How it is presented}

\begin{longtable}{@{}p{3.2cm}p{11.2cm}@{}}
\toprule
Option & Meaning \\
\midrule
\endhead
\texttt{--color} &
Force ANSI truecolour output, using 24-bit background colours and the standard
cubing palette. \\
\texttt{--no-color} &
Draw the net with coloured square characters instead, so the output survives being
redirected to a file or pasted into a document.  The default is automatic: colours
when standard output is a terminal, squares otherwise. \\
\texttt{--json} &
Print the machine-readable payload instead of the human listing. \\
\texttt{-h}, \texttt{--help} &
Usage summary. \\
\bottomrule
\end{longtable}

\subsection{Exit codes}

\begin{center}
\begin{tabular}{@{}cl@{}}
\toprule
Code & Meaning \\
\midrule
0 & A solution was produced (and verified). \\
2 & The input could not be understood --- bad notation, an impossible state --- or
the cube could not be solved.  The reason is written to standard error, prefixed
with \texttt{error:}. \\
\bottomrule
\end{tabular}
\end{center}

\subsection{Examples}

\begin{lstlisting}[style=shell]
# a scramble from a competition
uv run rubiks-cfop --scramble "D B2 F2 L2 U' F2 U B2 R2 U"

# reproducible random-state scramble, plain output for a log
uv run rubiks-cfop --random --seed 7 --no-color

# solve a state produced elsewhere, with a white cross of my choosing
uv run rubiks-cfop --facelets UULUUFUUFRRUBRRURRFFDFFUFFFDDRDDDDDDBLLLLLLLLBRRBBBBBB \
                   --cross white

# machine readable, for a script
uv run rubiks-cfop --scramble "R U R' U'" --json | jq .total_moves
\end{lstlisting}
